\documentclass[../main.tex]{subfiles}
\usepackage[utf8]{inputenc}
\usepackage[T1]{fontenc}
\usepackage{graphicx}
\usepackage{longtable}
\usepackage{wrapfig}
\usepackage{rotating}
\usepackage[normalem]{ulem}
\usepackage{amsmath}
\usepackage{amssymb}
\usepackage{capt-of}
\usepackage{hyperref}
\usepackage{float}
\graphicspath{{../}}
\author{Cezary Wieczorkowski}
\date{\today}
\title{Abstract}
\hypersetup{
 pdfauthor={Cezary Wieczorkowski},
 pdftitle={Abstract},
 pdfkeywords={},
 pdfsubject={},
 pdflang={Polish}}
\begin{document}

\selectlanguage{polish}
\begin{abstract}
Celem pracy było zbadanie możliwości realizacji wybranej splotowej sieci neuronowej (CNN) w systemie wbudowanym opartym o układ FPGA. W ramach pracy przeprowadzono trening sieci przy 
użyciu biblioteki Brevitas oraz implementacje sprzętową w układzie FPGA przy pomocy biblioteki Finn. Na końcu dokonano ewaluacji zaimplementowane sieci pod kątem wydajności obliczeniowej oraz energetycznej.
\end{abstract}
\textbf{Słowa kluczowe:} Technika cyfrowa, AI, sieci neuronowe, FPGA \\
\textbf{Dziedzina nauki i techniki, zgodne z wymogami OECD:} nauki inżynieryjno--techniczne: automatyka, elektronika, elektrotechnika i technologie kosmiczne
\newpage
\selectlanguage{english}
\begin{abstract}
The goal of this thesis was to study the possibilities of implementing a chosen convolutional neural network in an embedded system based around an FPGA. As part of this work a training of the neural network using the Brevitas
library and implementation in FPGA hardware using Finn library were conducted. Then the achieved implementation was evaluated in terms of compute and power efficiency.
\end{abstract}
\textbf{Keywords:} Digital electronics, AI, neural networks, FPGA \\
\selectlanguage{polish}
\end{document}
